\section{Jacobian rank and local rigidity}

\begin{theorem}[982-J01, 982-J02, 982-J03]
Let \(P\) be \(n\ge3\) distinct planar points with no three collinear. Let \(s\)
count pins with at least \(\lfloor n/2\rfloor\) distances, and let \(J\) be the real
Jacobian of all true equations \(|q-p|^2-|r-p|^2=0\) with distinct \(p,q,r\), evaluated
at \(P\). Write \(t=\dim_{\mathbb R}\ker J\). Then
\[
 \operatorname{trdeg}_{\mathbb Q}\mathbb Q(P)\le t\le 2s+576H_{n-s}.
\]
Here \(H_j=\sum_{i=1}^j i^{-1}\) and \(H_0=0\). By fixing \(t\) suitable scalar
coordinates to their actual values, the remaining coordinates are locally
unique. Over the field \(F\) generated by the fixed values,
\([F(P):F]\le 2^{2n-t}\).
\end{theorem}

\begin{proof}
Each row for a pin \(p\) and targets \(q,r\) has blocks
\[
 2(r-q),\qquad 2(q-p),\qquad 2(p-r),
\]
respectively. At each pin, order each class of targets and retain only
equalities between consecutive targets in that class. Two point blocks then
co-occur in at most six rows. There are at most two adjacencies when one is the
pin and the other a target, for either pin choice. When both are targets their
bisector contains at most two pins, and the target pair is adjacent at most
once at each of them. A fixed target's row block is collinear with its vector
to the pin. Each such direction occurs at most twice among its target roles:
one pin lies on that line and contributes at most two adjacencies.

The Dvir--Garg--Oliveira--Solymosi block rank theorem \cite{dgos} gives
corank at most \(2m/(1+k/24)\) for a matrix on \(m\) two-coordinate blocks
with at most three nonzero blocks in a row, overlap at most six rows between
two columns, and \(k\) well-spread \(1\times2\) blocks in every column. Here
well-spread means that no direction contributes more than \(k/2\), since the
only nontrivial test subspaces in dimension two are lines. All blocks are
real and the matrix rank is unchanged on passing to \(\mathbb C\).

In the velocity space \(\ker J\), freeze all \(s\) qualifying pins at a cost
of at most \(2s\) dimensions. A point equidistant to three frozen points has
velocity zero in the remaining space, since the two independent differences
of those point vectors are orthogonal to its velocity by the actual rows of
\(J\). Freeze such points for free as long as possible. If \(b\) points are
still unfrozen, every unfrozen pin sees at least \((n-b)/2\) distances into
the frozen targets. In the path ordering put all frozen targets first and
then retain only edges whose latter target is unfrozen. Orderings may be
chosen anew at this stage. Each unfrozen pin has at least
\[
 b-1-\bigl(d_P(p)-d_{\rm frozen}(p)\bigr)\ge b/2
\]
retained rows. This follows because every exclusively unfrozen class loses
only its first target from the count of latter targets.

For \(b\ge48\), put \(k=\lfloor b/12\rfloor\). Repeatedly remove an unfrozen
point currently occurring as a target fewer than \(k\) times, together with
its incident rows in any role. This cannot reach \(u=\lfloor b/3\rfloor\)
remaining points. The \(u\) remaining pins started with at least \(ub/2\)
rows and lost at most \((b-u)k\) to deletions: their centers were never
deleted. Their surviving rows number at most \(u(u-1)\) because the latter
target is a different remaining point and is unique at a pin. But for
\(b=3u+r\), \(0\le r\le2\), one has
\[
 12\bigl(u(b/2-u+1)-(b-u)b/12\bigr)=ru+12u-r^2>0.
\]
The surviving \(m>b/3\) points therefore each have \(k\ge4\) target-role
rows, with direction multiplicity at most two. Their velocity projection
lies in the kernel of the corresponding matrix on these \(m\) points,
since the other columns in the retained rows are frozen. Freezing this
batch costs at most \(2m/(1+k/24)\le576m/b\) dimensions. Summing these costs
using \(m/b\le H_b-H_{b-m}\), and charging the final \(b<48\) points directly
at cost at most \(2b\le576H_b\), proves the bound for \(t\).

Derivations with respect to a transcendence basis chosen from the actual
coordinates extend to the finitely generated coordinate field, since the
remaining extension is algebraic and characteristic zero. Applying them to
the equality equations yields independent tangent vectors: their entries
on basis coordinates are the identity. This proves the first inequality
without a smoothness assumption on the entire solution set.

Choose \(t\) coordinate functionals that give an injection on \(\ker J\).
The other \(2n-t\) columns of \(J\) are independent. A square submatrix on
these columns has nonzero determinant, so the corresponding equations after
fixing the \(t\) coordinates define an isolated simple root at \(P\).
The inverse function theorem gives local uniqueness. The derivation argument
over \(F\), applied to this invertible matrix, gives transcendence degree
zero over \(F\). The remaining coordinates are a finite separable extension.
Each embedding over \(F\) produces a different isolated simple complex
solution of the selected equations, all of degree at most two. The affine
B\'ezout bound for isolated solutions gives at most \(2^{2n-t}\) such
embeddings \cite{heintz}.
\end{proof}
